\subsection{Kurze Beschreibung zum Projektinhalt}
\label{subsec:Projektbeschreibung}

\par Das Programm \gls{r2h} ist mit einem Budget von 80 Mio. Euro und dem Ziel der Migration der 44 Rechnungskreise aus dem \gls{r3} ins \gls{s4} bis zum Jahr 2028 das bisher geplant umfangreichste Transformationsprojekte der Unternehmensgeschichte. Die Komplexität ergibt sich maßgeblich aus der über zwei jahrzehntelangen Historie des Altsystems und der zentralen Bedeutung des Hauptbuchs für die Bilanzierung und Steuerung des Konzerns. Um diese Risiken zu beherrschen, wurde ein iterativer Wellenansatz gewählt, der einen parallelen Betrieb von \gls{r3} und \gls{s4} während der Migrationsphase vorsieht. Die Einbindung zahlreicher SAP-Komponenten wie \gls{mbc} und \gls{aif} sowie die Neuintegration von Tochtergesellschaften wie der \gls{gls-adn} erhöhten die technischen und organisatorischen Abhängigkeiten. Diese Vielschichtigkeit der Stakeholderlandschaft – vom Konzernvorstand bis zu den Regionalclubs – erforderte eine stringente Koordination der Liefergegenstände über alle Projektphasen hinweg.


\begin{comment}
\par Ziel des Programms \ac{r2h} ist die konzernweite Einführung des neuen \ac{s4} und die Migration aller 44 im \ac{r3} befindlichen \ac{bk} (inkl. Schatten-\ac{bk}) bis zum 01.01.2028 (zuzüglich Hypercare) mit Gesamtprojektkosten von 80 Mio. Euro. Zur Erreichung dieses Ziels wurde ein iterativer Wellenansatz gewählt, sodass die \ac{bk} sukzessive in das \acs{s4} überführt werden. Ursprünglich in drei Wellen geplant, wurde deren Anzahl erhöht, da die erste Welle aufgrund zahlreicher Individualentwicklungen im Altsystem erhebliche Herausforderungen aufwies. Der Wellenansatz erfordert während der Koexistenz von \ac{r3} und \ac{s4} einen Parallelbetrieb beider Systeme.

\par Nach Programmstart wurden weitere SAP-Komponenten wie \acl{mbc} (\ac{mbc}), \acl{aif} (\ac{aif}), \acl{iag} (\ac{iag}), \acl{drc} (\ac{drc}) und \acl{refx} (\ac{refx}) ergänzt, um die Systemlandschaft des \ac{adac} zu modernisieren und den künftigen Betrieb zu vereinheitlichen. Zudem wurde die Integration des Luftfahrtunternehmens \acl{adn} (\ac{adn}) in das \ac{s4} beschlossen, welche im separaten Teilprojekt \acl{f2h} (\ac{f2h}) umgesetzt wird. Die branchenspezifischen Anforderungen und die dynamische IT-Landschaft der \ac{adn} erhöhten die Komplexität der \ac{ricefw}-Objekte und der Programmplanung.

\par In Welle 1 wurden die \ac{bk} der \ac{ais} (3800) und der \acl{avm} (\ac{avm}) (2400) ausgewählt. Der \ac{bk} 3800 wurde wegen seiner geringen Komplexität priorisiert; der \ac{bk} 2400 diente als Belastungstest für die Massentauglichkeit und Schnittstellenfähigkeit des \ac{s4}-Systems, insbesondere im Zusammenhang mit \ac{aif}. Die Produktivsetzung des \ac{bk}s 3800 erfolgte planmäßig am 01.01.2026. Für seine Implementierung wurden 150 \ac{ricefw}-Objekte umgesetzt, die die Anforderungen der \acl{ves} (\ac{ves}) im \ac{s4} abbilden.

\par Der Go-Live des \ac{bk}s 2400 erfolgte im Rahmen der Hauptwelle 1 in der Nebenwelle 1.1 aufgrund von Testkapazitätsengpässen und eines Stopps der Schnittstellenmodernisierung erst am 01.04.2026.\footnote{Die Bezeichnungen Hauptwelle, Nebenwelle und Patchwelle entsprechen der semantischen Versionsnummer.} Für 2026 sind neben ungeplanten Patches die Minorreleases Welle 1.2 mit \ac{drc} am 01.07. und Welle 1.3 \ac{f2h} am 01.10. sowie Implementierung und Abnahme der Welle 2 mit zwölf geplanten \ac{bk} am 01.01.2027 vorgesehen.

%To do Abkürzungen Implementieren.

\par Die Vielzahl interner und externer Stakeholder – darunter der Konzernvorstand, die 18 Regionalclubs, Fachbereiche wie Finanzen \ac{fnz}, Controlling, Einkauf, Recht und Compliance, die \ac{ais}, Datenschutz, diverse Tochtergesellschaften, Implementierungspartner sowie interne und externe Auditoren – verdeutlicht die Komplexität des Programms. Aufgrund der starken Abhängigkeiten zum Hauptbuch müssen zudem andere Projekte mit Schnittstellenbezug, etwa \ac{vita}, fortlaufend berücksichtigt werden.

\par Intensive Kommunikation war notwendig, um die betroffenen Gesellschaften und Mitarbeitenden über systemische Neuerungen und veränderte Geschäftsprozesse zu informieren und Akzeptanz zu schaffen.

\par Die Gesamtkomplexität des Programms und des Teilprojekts \ac{ricefw} ergibt sich aus mehreren Faktoren: der 25 Jahre alten \ac{r3}-Historie, dem zentralen Stellenwert des Hauptbuchs für Bilanzierung und Jahresabschluss, der großen und heterogenen Stakeholderlandschaft, unterschiedlichen Zielsetzungen, einer neuen technologischen Plattform sowie einem zuvor eher evolutionär arbeitenden internen Team.

\par Hinzu kam, dass viele Beteiligte erstmals in dieser Konstellation zusammenarbeiteten. Dies erforderte den Aufbau neuer Berichtswege, Abstimmungsformen und Mechanismen zur Einforderung von Ergebnissen gegenüber externen Partnern ohne direkte Disziplinargewalt.
\end{comment}